% New standalone increment. Prior sources and GitHub are read-only.
\documentclass[11pt]{article}
\usepackage[T1]{fontenc}\usepackage[utf8]{inputenc}\usepackage{lmodern}
\usepackage[margin=.88in]{geometry}
\usepackage{amsmath,amssymb,amsthm,mathtools,microtype,array,longtable,booktabs,xurl}
\usepackage[hidelinks]{hyperref}
\setlength{\emergencystretch}{3em}\allowdisplaybreaks[2]
\numberwithin{equation}{section}
\newtheorem{theorem}{Theorem}[section]
\newtheorem{lemma}[theorem]{Lemma}\newtheorem{proposition}[theorem]{Proposition}
\theoremstyle{remark}\newtheorem{remark}[theorem]{Remark}
\newcommand{\dd}{\,d}\newcommand{\E}{\mathbb E}
\newcommand{\KK}{\mathcal K}\newcommand{\BB}{\mathcal B}
\newcommand{\Gp}{{G_+}}\newcommand{\eps}{\varepsilon}
\newcommand{\Kfun}{\mathsf K}\newcommand{\Ffun}{\mathsf F}
\newcommand{\one}{\mathbf1}\newcommand{\ph}{\phi(0)}
\newcommand{\src}[2]{\textup{#1} \texttt{#2}}
\title{Strong-reference reduction to the selection error form\\
and the local initial guard}
\author{Standalone analytical increment for independent review}
\date{Version 1 --- 30 October 2026}
\begin{document}\maketitle
\begin{abstract}
On the initialized deep-regime passage, CFA's same-field reference
obeys TAILREP's strong error inequality with a pointwise additive
defect $O(q^{1/20})$ and a model-only coefficient of $NW_2$.
The selected weak distance uses the actual incoming probability
transported with fixed labels. Its conversion from current radial
weights is proved, and all non-selected weak moments are charged
with transverse energy rather than mass alone. The primary gate
remainder is superpolynomially small on the late interval. Before
that interval, CFA's original local inequality gives an initial
strong error $O(q^{1/20})$ without invoking the later local box.
The weak error inequality and selection remain separate obligations.
\end{abstract}

\section{Step 1: objects, intervals and hypothesis inventory}
\label{an03srrv1:inventory}
All new labels have prefix \texttt{an03srrv1:}. Fix the initialized
deep regime $\kappa=10$, $\rho\in[13/20,7/10]$, and write
$\lambda=\rho^2$, $Q=\lambda+q^2$, $\ell=\log(1/q)$,
$k=10(1-\lambda)$, and $a=a_\rho=\rho\Kfun(\rho a)$.
Keep RB's $t_0=t_h$ and $t_b=\inf\{c=c_b\}$. All estimates
below take place on RB's completed passage.
The fixed populations are $\KK,\Gp,\Sigma,\BB$, with
$S=\mathcal C_q\subset\Gp$, $P_q=\Gp\setminus S$ and
$\mathcal D=\varnothing$. Put $G=\Gp\cup\Sigma$,
$O=\Gp\cup\Sigma\cup\BB$, $M=\mu_\KK$ and
$\widetilde M=\min(M,1)$.

Use MP's start $t_s=\inf\{H=q^{1/8}\}$ and WBR's upper endpoint
$t_e=\inf\{H=h\}$. The exact logistic comparison is
\[
 N'=\lambda N(1-N),\qquad N(t_s)=N_0=q^{1/8}.
\]
For any later fixed choice $0<n<h$, let $t_n$ be its logistic hit.
MP places $t_n$ inside $(t_s,t_e)$ for small $q$. No choice of the
selection constants $C_w,L_*,n$ is made here.

There is one reference object throughout:
\begin{equation}
 z_c=(x_c,y_c)=(c_s,d_s),\quad
 e=|z_c-(a,a)|_2,\qquad E_s=|z_c-Z_*(N)|_2,
 \quad Z_*=(X_*,Y_*).
 \label{an03srrv1:reference}
\end{equation}
It is MOMC's actual characteristic, retained by CFA, not a mean,
not reset to $(a,a)$, and not evolved under $G_M$. Both errors are
Euclidean. The core law is the positive original radial measure,
normalized in CFA by
$\widetilde\mu_\KK=\mu_\KK\min(1,M^{-1})$.
Its normalized first moment about $z_c$ obeys
\[
 \bar D_\KK=M^{-1}\int_\KK(|x-x_c|+|y-y_c|)\dd\mu
       =O(q^{1/20}).
\]
The unnormalized $D_\KK$ is $M\bar D_\KK$; WEAK's $D_s$ is
the corresponding moment of its unit-mass strong family.
These normalizations are not interchanged without a mass defect.

For every $t\in[t_s,t_n]$ define the \emph{fixed incoming} probability
and distance
\begin{equation}
 \eta_s(\dd b)=\frac{m(t_s,b)}{N_S(t_s)}\dd\lambda_0(b),\quad
 W_2(t)=\big\||u(t,b)-U_*(N(t))|+|v(t,b)-V_*(N(t))|\big\|_{L^2(\eta_s)},
 \label{an03srrv1:W}
\end{equation}
where $N_S=\mu_S$, $u=(\pi/2-\theta)/q$ and
$v=(\psi-\pi/2)/q$. This is the convention S4e must retain.
The original labels, not their current radial weights, transport
$\eta_s$. For fixed $q$ the reached chart ensures finite moments.

\paragraph{Inventory of the consumer, not an application of it.}
TAILREP \texttt{selection} requires exact logistic $N$,
$0<N_0<n<1$, $N(t_n)=n$, nonnegative errors, and, up to the first
exit from $E_s<\delta$,
\begin{align}
 D^+E_s&\le-gE_s+C_s(d_0+NW_2),\notag\\
 D^+W_2&\le[-\lambda(1-N)/2+N/2+\ell_B(t)]W_2+C_w(E_s+d_0).
 \label{an03srrv1:consumer}
\end{align}
Here $g,C_s,C_w>0$ are fixed, $d_0\ge0$ is one pointwise
constant shared by both inequalities, $\ell_B\ge0$ and
$\int\ell_B\le L_*<\infty$. Weak return is maintained separately.
With
\begin{gather*}
 B_*=e^{L_*}(1-n)^{-1/(2\lambda)},\quad
 a_0=\lambda(1-n)/2,\quad m_0=\min(g/2,a_0/2),\\
 L_0=1+C_s(1+B_*)/(g-m_0),\quad K_0=2L_0,
 \quad\Theta=\frac{C_sB_*C_wn}{(a_0-m_0)(g-m_0)}\le1/4,
\end{gather*}
it also requires
\begin{equation}
 K_0(E_0+d_0+A_0\sqrt n)\le\delta/2,\qquad
 E_0=E_s(t_s),\quad A_0=\sqrt{N_0}W_2(t_s).
 \label{an03srrv1:consumerGuard}
\end{equation}
For terminal convergence $N_0,A_0,d_0\to0$ and common finite
constants are needed. This increment supplies only the first
inequality and the vanishing strong part of its initial guard.

WEAK's derivation of \texttt{strongineq, smallgain, guard} additionally
uses unit strong mass in its leading system, positive fixed label
measures, nonnegative returned weak inputs (also $U_*>0$), the
same logistic comparator and fixed normalized weak weights, finite
$W_p$, and a local box
\[
 \delta<a/8,\quad a+3\delta<2\ph,\quad
 |X_*-a|,|Y_*-a|\le\delta,
\]
whose frozen-weak-atom coherent Jacobian is at most $-g_WI$ on
the $2\delta$ box. Here
\[
 \alpha=\lambda\Phi(\rho a)/2,\quad
 \gamma=\frac{\alpha(1/2-\alpha)}{\alpha+1/2},\quad g_W=\gamma/2.
\]
Its strong moment error is pointwise, $D_s\le d_0$; the weak
return derivation uses $0\le Y\le a+3\delta<2\ph$,
$|\rho K_s|\le\rho\ph+\lambda(a+3\delta)$, and the strict
moment margin $C_un+U_0^+\sqrt{N_0n}<\ph-(a+3\delta)/2$.
Those leading assumptions do not become original-flow identities.
We instead use WBR's proved return, convert weights below, and
pay physical mass and moment defects. No term $\ell(t)E_s$ will
remain: every defect is pointwise additive or a model-only multiple
of $NW_2$. Thus no absorption based only on $\int\ell$ is used.

\section{Step 2: a pointwise strong-receiver remainder}
\label{an03srrv1:fieldsection}
CFA's proved representation on its fixed square
$[-L_c,L_c]^2$, $L_c=79/51+1$, is
\begin{align}
 2(F_c)_x&=-(1-\widetilde M)x+\rho\phi(\rho x)
 -\rho\widetilde{\mathcal S}(x)
 +\lambda[V_O+(1-c)y]\Phi(\rho x),\notag\\
 2(F_c)_y&=-(1-\widetilde M)y-\widetilde Y-(K_O)_+
                         +\rho(1-c)\Kfun(\rho x),
 \label{an03srrv1:cform}\\
 V_T&=-\frac{\E_{P_2}[(f_T)_1X_2]}{qQ},\qquad
 K_T=\frac{\E_{P_1}[(f_T)_2X_1]}q,
 \quad (V_O,K_O)=\sum_{T=S,P_q,\Sigma,\BB}(V_T,K_T).
 \label{an03srrv1:moments}
\end{align}
We start from its exact $F_{\rm orig}=F_c+\mathcal R$;
we do not re-derive the receiver field. Its full remainder
$\eps_F$ bounds the value needed here. To avoid an integrated
inference we restate and justify the pointwise majorant used in
CFA's proof.

\begin{lemma}[Primary gate tail, with a late rate]
\label{an03srrv1:tail}
With the fixed admissible $C_0$ already chosen for CFA's chart,
the primary selected-moment error has bound
\begin{equation}
 C_{L_c}\{J_\Gp+q\sqrt{E_\Gp J_\Gp}+\mu_\Gp\mathcal T(t)\},
 \quad
 \mathcal T(t)\le C\exp\left[-\frac{c_*}
 {q^2+C_*C_0^2e^{-b_*(t-t_0)}}\right]+Ce^{-c_*/q^2},
 \label{an03srrv1:tailpoint}
\end{equation}
where $b_*=Q(1-c_b)/2$. In particular, on $[t_s,t_e]$,
\begin{equation}
 \mathcal T(t)\le C\exp(-c q^{-597/320})+Ce^{-c/q^2}.
 \label{an03srrv1:taillate}
\end{equation}
It is smaller than every fixed power of $q$, with a constant
allowed to depend on that power. On $[t_0,t_b]$ it is simply
bounded by a fixed constant.
\end{lemma}
\begin{proof}
Restate the cancellation behind CFA \texttt{primaryselected}.
For $i=1,2$, write exactly
$f_{\Gp,i}=L_{i1}X_1+L_{i2}X_2+N_i$, with
$L_{ij}=\int_\Gp W_iU_j$ and
$N_i=\int_\Gp W_i(U\cdot X)_-$.
Subtracting the $P_2$ regression removes $L_{i2}X_2$.
The remaining linear terms cost $J_\Gp$ and
$q\sqrt{E_\Gp J_\Gp}$. For a unit input direction with
$u_2>0$, the scalar Gaussian negative-part moment is
\[
 \E_{P_2}(u\cdot X)_-^2
 \le Cq^2 e^{-c(u_2/q)^2}.
\]
Cauchy--Schwarz against the Gaussian factors, and
$|W_i||U|\le m$, bound the value errors after their possible
$1/q$ scaling by $C\mu_\Gp e^{-c(u_2/q)^2}$ uniformly over
the primary directions. RCET's ratio gives
\[
 q|U_1|/U_2\le C(C_0e^{-b_*(t-t_0)/2}+q^2),\quad
 (u_2/q)^{-2}\le C(q^2+C_0^2e^{-b_*(t-t_0)}).
\]
This proves the value part of \eqref{an03srrv1:tailpoint}.
For the full CFA norm, its receiver boundary has
$X_1=-qt_xX_2$, $|t_x|\le2L_c$.
The already fixed $C_0$ gives $U_2-qt_xU_1\ge U_2/2$,
so the negative part there vanishes for $X_2\ge0$.
The other side costs at most
$C\mu_\Gp q^{-1}\E[|X_2|^3\one_{X_2<0}]\le
C\mu_\Gp e^{-c/q^2}$. The regression correction is bounded
by the same value estimate times a bounded differentiated
Gaussian moment. This justifies the tail including both gate sides.

WBR's time formula, with its bounded $A_q$ shift, implies
\[
 e^{-b_*(t-t_0)}\le Cq^{(1-c_b)(k-1/8)/2}
                  \le Cq^{597/320},\qquad t\ge t_s.
\]
The last exponent uses $c_b\le1/4$ and $k\ge51/10$.
As $597/320<2$, the denominator in \eqref{an03srrv1:tailpoint}
is at most $Cq^{597/320}$. This proves \eqref{an03srrv1:taillate}.
The superpolynomial statement follows from the elementary boundedness
of $x^re^{-cx}$ for every fixed $r>0$, not from a numerical tail test.
\end{proof}

\begin{lemma}[Global moment charges and pointwise CFA remainder]
\label{an03srrv1:pointwise}
For an arbitrary cohort $T$,
\begin{equation}
 |V_T|+|K_T|\le C\left(\mu_T+q^{-1}\sqrt{\mu_TJ_T}\right),
 \qquad |V_\BB|+|K_\BB|\le C\mu_\BB/q.
 \label{an03srrv1:momentcharge}
\end{equation}
On the completed passage a valid pointwise majorant for CFA's
full remainder is
\begin{align}
 \eps_F\le C\big[&q^2(1+M_K+M_{\rm res}+\mu_G+|V_O|+|K_O|)
 +J_G+q\sqrt{E_GJ_G}+\mu_\Gp\mathcal T\notag\\
 &+E_\Sigma+q^{-1}\sqrt{E_\Sigma J_\Sigma}
 +\mu_\BB/q+(M-1)_++\mu_Gq^{-j}e^{-c/q^2}\big]
 \label{an03srrv1:majorant}
\end{align}
for fixed $j,C$. Consequently
\begin{equation}
 \sup_{[t_s,t_e]}\eps_F\le C_{\rm fixed}q^{1/10}.
 \label{an03srrv1:epslate}
\end{equation}
\end{lemma}
\begin{proof}
On $P_2$, $X=(qZ,V)$, balance gives
$|f_{T,1}|\le2qJ_T|Z|+2\sqrt{E_TJ_T}|V|$.
On $P_1$, $|f_{T,2}|\le2\sqrt{E_TJ_T}|X_1|+2E_T|X_2|$.
Insert these into \eqref{an03srrv1:moments}, use Gaussian moments
and $E_T\le\mu_T$ to prove \eqref{an03srrv1:momentcharge}.
The outside alternative follows directly from $|f_T(X)|\le\mu_T|X|$.

Here is the pointwise justification of every line of
\eqref{an03srrv1:majorant}, as opposed to its integrated use in CFA.
MOMC's compact core-kernel reduction, already used in CFA's exact
assembly, has error $C_{L_c,R}q^2(1+M_K)$; the core's second
output and response cost $Cq^2M_K$. Core mass capping costs
$C(M-1)_+$ in value (and receiver derivative), by subtracting
the two donor measures. CFA's $P_1$ Stein calculation gives
$J_G+q\sqrt{E_GJ_G}+q^2E_G+q^2|K_G|$, plus the fixed-chart
Gaussian receiver tail $C\mu_Gq^{-j}e^{-c/q^2}$.
Its $P_2$ trigonometric/kernel corrections give
$Cq^2(E_G+J_G+q^{-1}\sqrt{E_GJ_G})$ and
$Cq^2(|V_O|+|K_O|)$; the selected terms of $\Gp$ are exactly
Lemma \ref{an03srrv1:tail}. For $\Sigma$, inserting the two output
bounds above in the selected value integrals gives
$C(E_\Sigma+J_\Sigma+q^{-1}\sqrt{E_\Sigma J_\Sigma})$.
Boundary differentiation, when retaining the full norm, uses bounded
$t_x,t_x'$ and the same Gaussian moments, as in CFA; it has the
same charge. PROFILE's exact outside cost is $C\mu_\BB/q$,
also paying its formal moment replacements. Finally the primary
negative part of $K_O$ can occur only on $X_1<0$, since $W_2>0$
on $\Gp$; it costs $C\mu_\Gp q^{-1}e^{-c/q^2}$.
The secondary and outside negative parts are bounded by their
absolute charges above. This accounts for the positive-part
substitution and proves the majorant. No selected moment is
discarded by regression orthogonality.

RB gives $J_G\le q^2K_Y^2$, $E_G=O(1)$,
$\mu_\BB/q\le q^{1/10}$, and
$E_\Sigma\le C_\Sigma(q^5+q^{\Gamma_\nu}H)$ with
$\Gamma_\nu>3/10$. Thus the secondary line is
$O(q^{5/2}+q^{\Gamma_\nu/2})$ on the late interval.
The moments in the $q^2$ line are uniformly bounded by
\eqref{an03srrv1:momentcharge}. WBR gives
$(M-1)_+\le Cq^{11/10}$ on the whole passage, and the primary
tail is \eqref{an03srrv1:taillate}. These pointwise statements
prove \eqref{an03srrv1:epslate}. All constants multiplying vanishing
powers may use the already fixed passage bounds.
\end{proof}

The core first-moment replacement is a further, distinct value error:
\begin{equation}
 |\widetilde{\mathcal S}(x_c)-\widetilde M\Kfun(\rho x_c)|
 \le\rho\bar D_\KK,\qquad
 |\widetilde Y-\widetilde M y_c|\le\bar D_\KK.
 \label{an03srrv1:corecollapse}
\end{equation}
Indeed $\Ffun(z,w)$ is 1-Lipschitz in $w$ and
$\Ffun(z,z)=\Kfun(z)$; $\widetilde M\le1$.
This costs $O(q^{1/20})$ by RB/SI. The remaining replacement
$\widetilde M\mapsto1$ costs $C|1-M|=O(q^{11/10})$ on the
compact reference chart, by MP. No original characteristic is
replaced by the collapsed one in making these value comparisons.

\section{Steps 2(c)--3: weak moments and fixed incoming weights}
\label{an03srrv1:momentsection}
MP gives on $[t_s,t_e]$
\begin{gather}
 N_S/H=1+O(q^2),\quad
 d_{
m ph}:=\sup|\log(H/N)|
 \le C(q^{1/8}+q^2\ell+q^{1/4}h/b_0),\notag\\
 J_S\le C(q^3+q^2H),\quad
 \mu_{P_q}\le C(q^{1/4}H+q^2),\quad
 \mu_\Sigma\le C(q^5+q^{\Gamma_\nu}H+q^2),
 \label{an03srrv1:MPinputs}
\end{gather}
where $b_0=Q(1-c_b)$. The $q^2$ term in $\mu_{P_q}$ is retained:
MP's band statement is initially about $H_{P_q}$, not radial mass.

\begin{lemma}[Finite-noise selected moment values]
\label{an03srrv1:finiteweak}
With $L=C_L(1+q^{-603/640})$,
\begin{equation}
 \left|V_S-\int_S mv\right|+
 \left|K_S-\int_Sm\Kfun(u)\right|
 \le C_{\rm fixed}q^{37/320}.
 \label{an03srrv1:finitevalue}
\end{equation}
\end{lemma}
\begin{proof}
The reached chart gives $|qu|,|qv|\le1/2$ and
$\sin\theta=\cos(qu)\ge1/2$. The exact $P_2$ regression identity
therefore gives
\[
 V_S=\int_Sm\frac{\sin(qv)}q\cos(qu)
                      +O(\mu_Sq^{-1}e^{-c/q^2}).
\]
For $K_S$ there is a useful exact Gaussian formula. On $P_1$,
$(U\cdot X)/(q\sqrt m)=\sin(qu)/q+\sin(qu)Z_1+\cos(qu)Z_2$;
the centered Gaussian here has variance exactly one. Stein's identity
in $Z_1$ yields
\begin{equation}
 K_S=\int_Sm\cos(qv)\left[
 \Kfun\left(\frac{\sin(qu)}q\right)
 +q\sin(qu)\Phi\left(\frac{\sin(qu)}q\right)\right].
 \label{an03srrv1:Kexact}
\end{equation}
Use $|\Kfun'|\le1$, $\Kfun(u)\le\ph+|u|$, and the sine/cosine
Taylor bounds only on the reached chart. The two errors are at most
\[
 Cq^2(1+L^2)\int_Sm(1+|u|+|v|)
                  +C\mu_Sq^{-1}e^{-c/q^2}.
\]
Cauchy--Schwarz and MP's transverse theorem bound the integral by
$C(N_S+\sqrt{N_SJ_S}/q)\le C(H+\sqrt{qH})$.
This is uniformly bounded, and $q^2L^2=O(q^{74/640})$.
This proves \eqref{an03srrv1:finitevalue}, without applying a bounded
weak-chart reduction to the growing chart.
\end{proof}

\begin{lemma}[Value conversion of radial weights]
\label{an03srrv1:weights}
Let $\eta_t=m(t,b)\dd\lambda_0/N_S(t)$ on the fixed clock.
Then its density relative to \eqref{an03srrv1:W} satisfies
\begin{equation}
 \omega(t,b):=\frac{\dd\eta_t}{\dd\eta_s}(b),\qquad
 \sup_{t,b}|\omega-1|
 \le C[D_q+q^2(1+L^2)]=o(1),\quad
 D_q=C(q^{1/8}+q^2\ell).
 \label{an03srrv1:weightbound}
\end{equation}
This is a value bound only, not a derivative-weight budget.
\end{lemma}
\begin{proof}
BULK's individual and aggregate amplitude defects, equivalently its
\texttt{bulkweights} proof, give
\[
 e^{-2D_q}\le
 \frac{z_b^2(t)/H(t)}{z_b^2(t_s)/H(t_s)}\le e^{2D_q}.
\]
On the reached chart, per-label balance is
$m=z^2+a_-^2+j^2$, where $j^2=(U_1^2+W_1^2)/2$.
WBR alignment gives $a_-^2/m\le Cq^2$, while
$j^2/m\le Cq^2L^2$. Hence $m/z^2=1+O(q^2(1+L^2))$
uniformly, including both closed endpoints. Combining the two
times with MP's $N_S/H=1+O(q^2)$ proves
\eqref{an03srrv1:weightbound}. Positivity of $z$ and radial weights
justifies every division; the ratios extend to endpoint labels by
the same individual inequalities. Nothing differentiates $\omega$.
\end{proof}

\begin{proposition}[Weak moment comparison in the required norm]
\label{an03srrv1:weakmoments}
One may take the model constant $C_1=8$ so that
\begin{align}
 |c-N|&\le C_{\rm fixed}(q^{1/8}+q^2\ell),\notag\\
 |V_O-NV_*|+|K_O-N\Kfun(U_*)|
 &\le C_1NW_2+C_{\rm fixed}q^{1/10}.
 \label{an03srrv1:momentcomparison}
\end{align}
The same second bound holds with $K_O$ replaced by $(K_O)_+$.
\end{proposition}
\begin{proof}
The first line is MP's response and phase at $H\le h$.
For $S$, write the leading moment differences against the current
probability, then use \eqref{an03srrv1:weightbound} and
$N_S/N=1+O(d_{\rm ph}+q^2)$. For small $q$ these two ratios
are at most two. Since $\Kfun$ is 1-Lipschitz,
\begin{align*}
 &\left|\int_S mv-NV_*\right|+
   \left|\int_Sm\Kfun(u)-N\Kfun(U_*)\right|\\
 &\hspace{1cm}\le4N\int_S(|v-V_*|+|u-U_*|)\dd\eta_s
                    +C N(d_{\rm ph}+q^2)\\
 &\hspace{1cm}\le4NW_2+C N(d_{\rm ph}+q^2).
\end{align*}
Here AN06 gives $0<U_*<1$ and $0<V_*<11/10$ for $N<2h<1/2$.
This comparison avoids multiplying a small weight error by the
possibly divergent supremum $L$; deviations are kept in $NW_2$.
Add Lemma \ref{an03srrv1:finiteweak}.

For the other primary labels, \eqref{an03srrv1:momentcharge} and
\eqref{an03srrv1:MPinputs}, with $J_{P_q}\le q^2K_Y^2$, give
\[
 |V_{P_q}|+|K_{P_q}|\le
 C_{\rm fixed}(q^{1/4}H+q^2+q^{1/8}\sqrt H+q)=O(q^{1/8}).
\]
For $\Sigma$ the analogous bound is
$C_{\rm fixed}(q+q^{\Gamma_\nu/2})$, and for $\BB$ it is
$C\mu_\BB/q\le Cq^{1/10}$. These are field-moment charges,
not division of MP's relative-mass estimate by $q$.
All remaining powers, including $37/320$ and $\Gamma_\nu/2$,
exceed $1/10$. This proves the second line, with slack $C_1=8$.
Finally $N\Kfun(U_*)\ge0$ and the positive-part map is
1-Lipschitz; its replacement needs no separate negative-part
estimate in this moment comparison.
\end{proof}

\section{Step 4: a model-only local box and the strong inequality}
\label{an03srrv1:local}
The coherent strong vector field with selected weak atom held fixed is
\begin{equation}
 F_*(x,y;N)=\frac12\binom{
 \lambda[NV_*-x+(1-N)y]\Phi(\rho x)}{
 -y-N\Kfun(U_*)+\rho(1-N)\Kfun(\rho x)}.
 \label{an03srrv1:coherent}
\end{equation}
Its symmetric Jacobian is exactly WEAK's displayed matrix
\begin{equation}
 J_*=
 \begin{pmatrix}
 \frac\lambda2[-\Phi(\rho x)+\rho(NV_*-x+(1-N)y)\phi(\rho x)]
 &\frac\lambda2(1-N)\Phi(\rho x)\\
 \frac\lambda2(1-N)\Phi(\rho x)&-1/2
 \end{pmatrix}.
 \label{an03srrv1:Jacobian}
\end{equation}
At $(a,a,0)$ it is $J_1\le-\gamma I$; LOC/CFA give the
model lower bound $\gamma\ge\gamma_*:=1183/20800$.

Here are explicit analytic choices, without a numerical sign test.
Let $L_J\ge1$ bound the sum of the operator norms of the first
derivatives in $x,y,N$ of \eqref{an03srrv1:Jacobian}, with $V_*$
treated as an independent frozen parameter in $[0,11/10]$, on
$|x-a|,|y-a|\le1/4$, $0\le N\le1/2$, over the ratio box.
The formula and the bounded Gaussian derivatives prove $L_J<\infty$.
Its value at $N=0$ is independent of that frozen parameter, so
\[
 \|J_*-J_1\|\le L_J(|x-a|+|y-a|+N).
\]
Use the model bounds $a_{\min}=39/160\le a< a_{\max}=2/5$ and choose
\begin{equation}
 \delta=\frac12\min\left\{
 \frac{a_{\min}}8,\frac{2\ph-a_{\max}}3,\frac18,
 \frac{\gamma_*}{8L_J}\right\}>0.
 \label{an03srrv1:delta}
\end{equation}
P \texttt{pa:branch} gives a $C^1$ graph at $N=0$.
There are model constants $n_{\rm br}>0,C_{\rm br}<\infty$ such that
\begin{equation}
 |Z_*(N)-(a,a)|_2\le C_{\rm br}N\qquad(0\le N\le n_{\rm br}).
 \label{an03srrv1:branch}
\end{equation}
For uniformity, the resident equations and the vector field are smooth
in the compact ratio parameter, the unstable eigenvalue is
$\lambda\ge169/400$, and the angular eigenvalues are separated
from zero by the model bounds used in P/LOC/AN06. The parameter
version of the same finite-dimensional unstable-manifold theorem
used in P gives a uniform $C^1$ graph; its first derivative on a
smaller compact graph gives $C_{\rm br}$. Equivalently one may use
P/AN06's first-order branch expansion with this compactness check.

Decrease the fixed WBR level, if necessary, so that
\begin{equation}
 h\le\min\left\{h_{\rm WBR},\frac{n_{\rm br}}2,
       \frac{\delta}{2C_{\rm br}},\frac{\gamma_*}{8L_J}\right\},
 \qquad h<1/4.
 \label{an03srrv1:hbound}
\end{equation}
The symbol $h_{\rm WBR}$ denotes any fixed number satisfying all
WBR's strict upper conditions. Then both local requirements hold:
$|X_*-a|,|Y_*-a|\le\delta$ for $N\le2h$, and, on the $2\delta$ box,
\[
 J_*\le-\gamma I+L_J(4\delta+2h)I
       \le-(\gamma/2)I=-g_WI.
\]
These choices also give $\delta<a/8$ and $a+3\delta<2\ph$.
They do not presume that the actual reference has entered this box.

\begin{proposition}[Strong error reduction up to its local first exit]
\label{an03srrv1:strong}
There is a pointwise constant
\begin{equation}
 d_0^{(s)}=C_{\rm srr}q^{1/20}\longrightarrow0
 \label{an03srrv1:d0}
\end{equation}
such that, on $[t_s,t_n]$ up to the first exit from $E_s<\delta$,
\begin{equation}
 D^+E_s\le-g_WE_s+C_s(d_0^{(s)}+NW_2),\qquad C_s=980.
 \label{an03srrv1:strongerror}
\end{equation}
$C_s$ and $g_W$ are model-only; $C_{\rm srr}$ may depend on
the already fixed passage constants. No $\ell(t)E_s$ term is left.
\end{proposition}
\begin{proof}
At $z_c$, use \eqref{an03srrv1:corecollapse} in CFA's field and then
replace $\widetilde M$ by one. On the local box $|x_c|,|y_c|<1$,
the coefficients of the moment and mass differences are bounded by
model constants. Directly from \eqref{an03srrv1:cform}, one may use
the conservative value bound
\begin{align*}
 |F_{\rm orig}(z_c)-F_*(z_c;N)|_2
 \le{}&10\{\bar D_\KK+|1-M|+|c-N|
       +|V_O-NV_*|+|(K_O)_+-N\Kfun(U_*)|\}+\eps_F.
\end{align*}
For example $\lambda\Phi/2\le1/2$, the coefficient of $K_O$
is $1/2$, and the remaining coefficients involve only $|x_c|,|y_c|<1$
and bounded $\Kfun(\rho x_c)$. Equations
\eqref{an03srrv1:epslate}, \eqref{an03srrv1:momentcomparison} and
MP's mass bound make this at most
$80NW_2+C_{\rm fixed}q^{1/20}$.
Choose $C_{\rm srr}$ to dominate the additive constant.
WEAK permits $C_s=20(R_*+1)^2$ with
$R_*>2+a+|u_\rho|+a/\lambda$; the model choice $R_*=6$ is
admissible by AN06's resident bounds. Thus $C_s=980$ dominates
both WEAK's prescribed coefficient and $80$.

The selected branch at the exact logistic mass satisfies
$Z_*'=F_*(Z_*;N)$. While $E_s<\delta$, the segment from
$z_c$ to $Z_*$ stays in the $2\delta$ box. The symmetric Jacobian
bound therefore gives
$(z_c-Z_*)\cdot(F_*(z_c;N)-F_*(Z_*;N))\le-g_WE_s^2$.
The norm Dini derivative gives \eqref{an03srrv1:strongerror},
also at $E_s=0$. All mass errors were charged as bounded values
on this box, and all weak spread was retained as $NW_2$.
No integrated error was substituted for $d_0^{(s)}$.
\end{proof}

\section{Step 5: the initial guard from the earlier CFA inequality}
\label{an03srrv1:entrysection}
This section uses neither WEAK's box nor $g_W$ before $t_s$.
The only attraction inequality is CFA \texttt{localineq}, already
certified by RB with $e\le3e_*/4<e_*$ and $r_\KK<1$:
\begin{equation}
 D^+e\le-g_c e+C_{\rm pass}(r_\KK+f_*),\quad
 g_c=\gamma_*/2,\quad
 f_*=c+|V_O|+(K_O)_++\eps_F.
 \label{an03srrv1:earlier}
\end{equation}
The capped donor mass stays in its proved range
$\widetilde M\in[1/2-o(1),1]$ and $(a,a)$ is stationary for
the comparison $G_{\widetilde M}$ throughout that range.

\begin{lemma}[Pointwise pre-start forcing]
\label{an03srrv1:preforcing}
On the entire interval $[t_0,t_s]$,
\begin{equation}
 c\le C(H+q^2+\mu_\BB),\qquad
 |V_O|+|K_O|\le C_{\rm fixed}q^{1/16},\qquad
 \eps_F\le C_{\rm fixed}q^{1/10},\qquad
 f_*\le C_{\rm fixed}q^{1/16}.
 \label{an03srrv1:prebounds}
\end{equation}
These statements do not use MP's late response or reached chart.
\end{lemma}
\begin{proof}
RB's \texttt{primaryclock} applies on all of $[t_0,t_b]$, not just
after $t_s$: its proof uses Q2EC's original entry seed and RCET's
fixed-subset transport under RB's residual/range/learning bounds.
With LCRC signgood it gives $E_\Gp\le CH$ before $t_s$.
The first-hit definition gives $H\le q^{1/8}$ there.
The secondary bound and the transverse guard consequently yield
$E_G\le C(H+q^5)$, $J_G\le q^2K_Y^2$.
LCRC's absolute charge, including Young's inequality, gives
$|c_G|\le2E_G+q^2J_G/Q$. Add the core $O(q^2)$ response and
$|c_\BB|\le C\mu_\BB$ to obtain the first bound.
The $C\mu_\BB$ term is retained: the outside supremum alone
does not justify deleting it in favor of $Cq^2$.

The proof of \eqref{an03srrv1:momentcharge}, in its sharper form
$C(E_G+J_G+q^{-1}\sqrt{E_GJ_G})$, and the outside alternative give
\[
 |V_O|+|K_O|\le C_{\rm fixed}
  [H+q^2+q^{5/2}+\sqrt H+\mu_\BB/q]
                  \le C_{\rm fixed}q^{1/16}.
\]
This is a balance/Gaussian estimate for arbitrary donors, not a late
weak-coordinate expansion. In \eqref{an03srrv1:majorant}, use
the bounded, not the late superpolynomial, estimate $\mathcal T\le C$.
It gives $\mu_\Gp\mathcal T\le C(H+q^2)$.
All transverse and finite-noise terms are $O(q^2)$, the secondary
term is at most
$C(q^{5/2}+q^{(\Gamma_\nu+1/8)/2})$, the outside term is
$q^{1/10}$, and WBR's global overshoot bound gives
$(M-1)_+\le Cq^{11/10}$. Thus $\eps_F\le Cq^{1/10}$.
Finally $\mu_\BB\le q^{11/10}$ and $H\le q^{1/8}$ bound
$c$ by $Cq^{1/8}$; summing proves the last assertion.
\end{proof}

\begin{proposition}[Vanishing local initial error]
\label{an03srrv1:E0}
For the same actual reference,
\begin{equation}
 e(t_s)=O(q^{1/20}),\qquad E_0=E_s(t_s)=O(q^{1/20})\to0.
 \label{an03srrv1:entryrate}
\end{equation}
\end{proposition}
\begin{proof}
Use \eqref{an03srrv1:earlier} only in its certified CFA regime:
\[
 e(t_s)\le e^{-g_c(t_s-t_0)}e(t_0)
 +C_{\rm pass}\int_{t_0}^{t_s}e^{-g_c(t_s-s)}
                      (r_\KK+f_*)(s)\dd s.
\]
RB has $r_\KK=O(q^{1/20})$ on this whole interval; Lemma
\ref{an03srrv1:preforcing} supplies the pointwise forcing.
The integral is at most $C(q^{1/20}+q^{1/16})/g_c$.
MOMC/CFA gives $e(t_0)=o(1)$, in particular bounded, and WBR's
time formula gives
\[
 e^{-g_c(t_s-t_0)}\le Cq^{g_c(k-1/8)/Q},\qquad
 \frac{g_c(k-1/8)}Q\ge\frac{235417}{832000}>\frac1{20}
\]
for small $q$, using $Q\le1/2$.
This proves the first bound. Then \eqref{an03srrv1:branch} at
$N_0=q^{1/8}$ gives
$E_0\le e(t_s)+C_{\rm br}q^{1/8}$.
In particular the required small forcing on a final stretch holds:
one may take $T_q=(t_s-t_0)/2\to\infty$; the stronger pointwise
bound was proved on the entire earlier interval. No logarithmic
length is multiplied by a fixed integrated forcing allowance.
\end{proof}

\section{Step 6: strong-error theorem and S4e interface}
\begin{theorem}[Original strong-reference part of the selection system]
\label{an03srrv1:main}
In the initialized deep regime, fix constants in the order
\[
 \mathrm{WPT/RB}\ \longrightarrow\ \mathrm{WBR}\
 \longrightarrow\ (a_s=1/8,p=2)\
 \longrightarrow\ (\delta,g_W,C_s)\
 \longrightarrow\ h\
 \longrightarrow\ [\mathrm{S4e}:C_w,L_*,n]\
 \longrightarrow\ q,
\]
where the model-only local choices are
\eqref{an03srrv1:delta}--\eqref{an03srrv1:hbound},
$g_W=\gamma/2$ and $C_s=980$.
For any fixed $0<n<h$, on $[t_s,t_n]$ up to the first exit from
$E_s<\delta$, CFA's actual reference and the fixed-incoming-law
$W_2$ in \eqref{an03srrv1:W} satisfy
\[
 D^+E_s\le-g_WE_s+C_s(d_0^{(s)}+NW_2),\qquad
 d_0^{(s)}=C_{\rm srr}q^{1/20}.
\]
Moreover $E_0=O(q^{1/20})$ and MP gives $A_0=O(q^{1/16})$.
For every fixed finite $K_0$ and fixed $n$, therefore,
\begin{equation}
 K_0(E_0+d_0^{(s)}+A_0\sqrt n)\le\delta/2
 \label{an03srrv1:strongguard}
\end{equation}
once $q$ is small. No extra retained inequality is assumed to obtain
this strong reduction or its initial error.
\end{theorem}
\begin{proof}
The field and moment estimates prove Proposition
\ref{an03srrv1:strong}; the independently justified earlier-regime
argument proves Proposition \ref{an03srrv1:E0}. MP's incoming
smallness uses the same probability and comparator convention as
\eqref{an03srrv1:W}. This proves \eqref{an03srrv1:strongguard}.
All original estimates used are confined to $[t_0,t_n]\subset[t_0,t_b]$.
\end{proof}

\paragraph{What this does not conclude.}
TAILREP ultimately needs $d_0=\max\{d_0^{(s)},d_0^{(w)}\}$.
If S4e proves $d_0^{(w)}\to0$ and fixed finite selection constants,
the same limit proves its full initial guard. This sentence does
not supply that missing weak estimate. The derivative of normalized
weights, the complete $\ell_B$ budget, $C_w,L_*$, the small-gain
choice $\Theta\le1/4$, the choice of $n$, and the selection theorem
remain in S4e. No strong guard is propagated to $t_n$ by this
increment alone; its differential inequality is stated with the
correct local first-exit scope. The pointwise $d_0^{(s)}$ is not
CFA's integrated $C_0E_g$ allowance.

\begin{thebibliography}{15}\small
\bibitem{RB} RB: \path{AN03_response_barrier_v1.tex},
\texttt{an03rbv1:completed, primaryclock, conclusions, massconclusions}.
\bibitem{CFA} CFA: \path{AN03_c_form_assembly_v1.tex},
\texttt{an03cfav1:representation, selected, primaryselected, cform,
majorant, localineq, passage}.
\bibitem{CLOCK} CLOCK: \path{AN03_exact_amplitude_clock_and_c_passage_v1.tex},
\texttt{inc:AN03:clock:v1:cfield, passage}.
\bibitem{LOC} LOC: \path{AN03_local_rate_mean_tracking_v1.tex},
local donor cancellation and Lemma 3.1.
\bibitem{WBR} WBR: \path{AN03_weak_boundary_return_v1.tex},
\texttt{an03wbrv1:main, mass, massrate, times, latechart}.
\bibitem{MP} MP: \path{AN03_matching_package_v1.tex},
\texttt{an03mpv1:accounting, phase, masses, transverse, smallness,
massbudgets}.
\bibitem{BULK} BULK: \path{AN03_bulk_clock_and_secondary_Q2_energy_v1.tex},
\texttt{inc:AN03:bulkclock:v1:bulkclockclosure, bulkweights}.
\bibitem{RCET} RCET: \path{AN03_residual_closure_and_eta0_transport_v1.tex},
\texttt{an03rcetv1:ratio, margin, transport}.
\bibitem{SI} SI: \path{AN03_selection_interface_v1.tex},
\texttt{an03siv1:inventory, coremoments, massdeficit}.
\bibitem{WEAK} WEAK: \path{AN03_weak_transit_forcing_and_orbit_selection_v1.tex},
\texttt{inc:AN03:weak:v1:strongineq, smallgain, guard}.
\bibitem{TAILREP} TAILREP: \path{AN03_mass_charged_tail_and_selection_defects_v1.tex},
\texttt{inc:AN03:tailrepair:v1:strongerror, weakerror, selection}.
\bibitem{P} P: \path{THEOREM_A_ANALYTICAL_PROGRESS.tex},
\texttt{pa:branch, pa:pop-x, pa:pop-y}.
\bibitem{AN06} AN06: \path{AN06_learned_orbit_witnesses_v2.tex},
\texttt{inc:AN06:v2:residentbounds, regiontheorem}.
\bibitem{PROFILE} PROFILE: \path{AN02_strong_entry_profile_and_tail_budget_v1.tex},
Proposition 4.1; exact outside value cost.
\bibitem{MOMC} MOMC: \path{AN02_moving_overlap_M_clock_passage_v1.tex},
strong compact-kernel reduction and actual-reference entry.
\end{thebibliography}
\end{document}
